\chapterhead{60}{ORR-21}{Capital Regulation Before the Global Financial Crisis}{1}{2}

\lohead{60}{a}{Explain the motivations for introducing the Basel regulations,
including key risk exposures addressed, and explain the reasons for revisions to
Basel regulations over time.}

\orsub{Early banking practices}

\begin{lettered}
\item Banks used substantial capital and grand buildings to establish trust, and
formed consortiums for self-regulation to prevent bank runs.
\item But this was inadequate during financial panics, necessitating
government-backed entities.
\item Increased international trade in the 1960s and 1970s required coordinated
financial regulations, and interconnected global banks posed risks to multiple
economies.
\end{lettered}

\orsub{Establishment of BCBS}

\begin{lettered}
\item Central bankers from the G10 nations, focusing on standardised global
regulations, formed the BCBS in 1974.
\item Basel I was introduced in 1988 to better gauge capital adequacy and
established minimum capital requirements based on risk exposures.
\end{lettered}

\orsubb{Capital types in Basel I}

\begin{checks}
\item \textbf{Tier 1 (core capital):} equity and specific preferred stocks for
solvency.
\item \textbf{Tier 2 (supplementary capital):} additional securities for loss
absorption.
\end{checks}

% ---------------------------------------------------------------------
\lohead{60}{b}{Explain the calculation of risk-weighted assets and the capital
requirement per the original Basel I guidelines.}

\orsub{Capital and risk-weighted assets (credit risk)}

Basel I required consolidated banking organisations to maintain:

\begin{orkeybox}
\[
\frac{\text{Tier 1 capital}}{\text{RWA}} > 4\%
\qquad\qquad
\frac{\text{Total capital}}{\text{RWA}} > 8\%
\]
\end{orkeybox}

\begin{center}
\begin{tikzpicture}[font=\fontsize{7.8}{9.6}\selectfont,
  hd/.style={text width=7.2cm,align=center,rounded corners=3pt,inner sep=5pt,
             text=white,font=\sffamily\bfseries\fontsize{8.4}{10}\selectfont},
  bd/.style={text width=7.2cm,align=left,rounded corners=3pt,inner sep=7pt,
             draw=palenavy,line width=0.6pt,text=navy}]
\node[hd,fill=navy,anchor=north west]  (ah) at (-7.9,0) {Assets};
\node[hd,fill=teal,anchor=north west]  (lh) at ( 0.4,0) {Liabilities};
\node[bd,fill=paleblue,anchor=north west] (ab) at (-7.9,-0.58) {
  $\blacktriangleright$~Cash / gold, claims on OECD governments\\[2pt]
  $\blacktriangleright$~Claims on OECD banks and on OECD public sector entities\\[2pt]
  $\blacktriangleright$~Uninsured residential mortgages\\[2pt]
  $\blacktriangleright$~Loans to consumers / commercial};
\node[bd,fill=paleblue,anchor=north west] (lb) at ( 0.4,-0.58) {
  \textbf{Deposits}\\[4pt]
  \textbf{Equity}\\[2pt]
  \emph{Tier 1 capital}\\
  \hspace{6pt}Equity (subtract goodwill from equity)\\
  \hspace{6pt}Noncumulative perpetual preferred stock\\[3pt]
  \emph{Tier 2 capital}\\
  \hspace{6pt}Cumulative perpetual preferred stock\\
  \hspace{6pt}99-year debentures\\
  \hspace{6pt}Subordinated debt $>5$ years};
\begin{scope}[on background layer]
  \node[fit=(ab)(lb),inner sep=0pt] (row) {};
\end{scope}
\node[hd,fill=rust,anchor=north west] (oh) at (-7.9,-4.55) {Off-balance sheet assets};
\node[bd,fill=palerust,anchor=north west] at (-7.9,-5.13) {
  $\blacktriangleright$~Guarantees (loans / bonds), bankers acceptances\\[2pt]
  $\blacktriangleright$~Letters of credit\\[2pt]
  $\blacktriangleright$~Interest rate swaps (current exposure)};
\end{tikzpicture}
\end{center}
\orfigcap{The Basel I balance sheet, on and off.}

\orsub{Risk-weighted assets (RWA)}

\begin{lettered}
\item RWA is a bank's assets weighted according to risk. This sort of asset
calculation is used in determining the capital requirement, or Capital Adequacy
Ratio, for a financial institution.
\item Assets are multiplied by their respective risk weights.
\end{lettered}

\begin{ordefbox}{Credit Equivalent Amounts (CEAs)}
Off-balance sheet items are converted into on-balance sheet equivalents using
prescribed conversion factors.
\end{ordefbox}

\noindent\begin{minipage}{\textwidth}
\renewcommand{\arraystretch}{1.32}
\begin{tabularx}{\textwidth}{@{}L{0.17\textwidth}Y@{}}
\rowcolor{navy}\multicolumn{2}{@{}l@{}}{\hspace{4pt}\thd{Figure 60.1: Risk weights for on-balance sheet items}}\\
\rowcolor{palenavy} \tsub{Risk weight} & \tsub{Asset category}\\
{\tabfont \textbf{0\%}} & {\tabfont Cash, gold, claims on Organization of Economic Co-operation and Development (OECD) countries such as U.S. Treasury bonds, and insured residential mortgages}\\
\rowcolor{paleblue}
{\tabfont \textbf{20\%}} & {\tabfont Claims on OECD banks and government agencies like U.S. agency securities such as Fannie Mae and Freddie Mac, or municipal bonds}\\
{\tabfont \textbf{50\%}} & {\tabfont Uninsured residential mortgages}\\
\rowcolor{paleblue}
{\tabfont \textbf{100\%}} & {\tabfont Loans to corporations, consumer loans, corporate bonds, claims on non-OECD banks}\\
\end{tabularx}
\end{minipage}
\orfigcap{Basel I on-balance sheet risk weights.}

\noindent\begin{minipage}{\textwidth}
\renewcommand{\arraystretch}{1.32}
\begin{tabularx}{\textwidth}{@{}L{0.20\textwidth}Y@{}}
\rowcolor{navy}\multicolumn{2}{@{}l@{}}{\hspace{4pt}\thd{Figure 60.2: Credit conversion factors}}\\
\rowcolor{palenavy} \tsub{Conversion factor} & \tsub{Off-balance sheet category}\\
{\tabfont \textbf{100\%}} & {\tabfont Guarantees on loans and bonds, bankers acceptances, and equivalents}\\
\rowcolor{paleblue}
{\tabfont \textbf{50\%}} & {\tabfont Standby letters of credit and warranties related to transactions}\\
{\tabfont \textbf{20\%}} & {\tabfont Loan commitments with original maturities equal to or greater than one year}\\
\rowcolor{paleblue}
{\tabfont \textbf{0\%}} & {\tabfont Loan commitments with original maturity less than one year}\\
\end{tabularx}
\end{minipage}
\orfigcap{Basel I credit conversion factors for off-balance sheet items.}

\begin{orexambox}{Example: risk-weighted assets}
The assets of Blue Star Bank consist of \$20 million in U.S. Treasury bills, \$20
million in insured mortgages, \$50 million in uninsured mortgages, and \$150
million in corporate loans. Using the risk weights from Figure 60.1,
\textbf{calculate} the bank's risk-weighted assets.

\vspace{6pt}
\noindent{\sffamily\bfseries\fontsize{8.6}{10}\selectfont\color{forest}Solution}\par
\[
(0.0\times\$20)+(0.0\times\$20)+(0.5\times\$50)+(1.0\times\$150)=\$175\text{ million}
\]
\end{orexambox}

\begin{ornotebox}{Source note}
The crash course references a table of add-on factors as Figure 60.3 but does not
reproduce it. The tables below are restored from the GARP curriculum reading,
\emph{Capital Regulation Before the Global Financial Crisis}, Chapter 21, which is
the source the crash course is condensing. They are not in your original PDF.
\end{ornotebox}

\orsub{Add-on factors ($D$) for derivative contracts}

Basel I offered authorities in each nation a choice between two methods of
computing a credit equivalent amount. This structure was revised in 1995 with the
addition of a maturity bucket greater than five years.

\orsubb{1) Current Exposure Method}

First, calculate the current market value of the contract, $V$. If the current
market value is negative, making it a liability rather than an asset, set $V=0$.
Second, add an amount $D$ to account for changes in the contract's future market
value.

\noindent\begin{minipage}{\textwidth}
\renewcommand{\arraystretch}{1.35}
\begin{tabularx}{\textwidth}{@{}L{0.34\textwidth}YY@{}}
\rowcolor{navy} \thd{Remaining maturity} & \thd{Interest rate swaps} & \thd{Foreign exchange swaps}\\
{\tabfont Less than one year} & {\tabfont zero} & {\tabfont 1.0\%}\\
\rowcolor{paleblue}
{\tabfont Five years or less} & {\tabfont 0.5\%} & {\tabfont 5.0\%}\\
{\tabfont More than five years} & {\tabfont 1.5\%} & {\tabfont 7.5\%}\\
\end{tabularx}
\end{minipage}
\orfigcap{Add-on factor $D$ as a percentage of notional, Current Exposure Method.}

\begin{orexambox}{Reading the table against the worked example}
The \$175 million interest rate swap below has three years remaining, which falls
in the ``five years or less'' band, giving $D=0.5\%$. That is the 0.005 used in
the solution.
\end{orexambox}

\orsubb{2) Original Exposure Method}

Available only for interest rate and foreign exchange contracts. Nations could
ignore the current market value of the contract and choose whether to use the
original or the remaining maturity.

\noindent\begin{minipage}{\textwidth}
\renewcommand{\arraystretch}{1.35}
\begin{tabularx}{\textwidth}{@{}L{0.34\textwidth}YY@{}}
\rowcolor{navy} \thd{Maturity} & \thd{Interest rate contracts} & \thd{Foreign exchange contracts}\\
{\tabfont Less than one year} & {\tabfont 0.5\%} & {\tabfont 2\%}\\
\rowcolor{paleblue}
{\tabfont Between one and two years} & {\tabfont 1\%} & {\tabfont 5\%}\\
{\tabfont Greater than two years} & {\tabfont $1\%+1\%\times \mathrm{INT}[M-1]$} & {\tabfont $5\%+3\%\times \mathrm{INT}[M-1]$}\\
\end{tabularx}
\end{minipage}
\orfigcap{Add-on factor $D$, Original Exposure Method. $\mathrm{INT}[X]$ returns the
closest integer to $X$.}

\begin{ornotebox}{Equity and commodity derivatives}
Equity and commodity derivatives were not discussed in original Basel I. The risk
weight was set according to the nature of the counterparty, except that no risk
weight could be more than 50\%. The 1995 Amendment introduced add-on factors for
such derivatives. The curriculum's own worked example applies 10\% to both equity
options and commodity (wheat) options.
\end{ornotebox}

\begin{orexambox}{Example: credit equivalent amount}
A bank has entered a \$175 million interest rate swap with a remaining maturity of
three years. The current value of the swap is \$2.5 million. Using the add-on
factors in Figure 60.3, \textbf{calculate} the swap's credit equivalent amount.
Also, calculate the risk-weighted assets assuming (1) the counterparty is an OECD
bank and (2) the counterparty is a corporation.

\vspace{6pt}
\noindent{\sffamily\bfseries\fontsize{8.6}{10}\selectfont\color{forest}Solution}\par
\[
\text{credit equivalent amount}=\$2.5+(0.005\times\$175)=\$3.375\text{ million}
\]
RWA assuming the counterparty is an OECD bank:
$\$3{,}375{,}000\times0.2=\$675{,}000$.\\
RWA assuming the counterparty is a corporation:
$\$3{,}375{,}000\times1.0=\$3{,}375{,}000$.
\end{orexambox}

\begin{orexambox}{Example: required capital}
Using the information from the previous examples, \textbf{calculate} Blue Star
Bank's required capital, assuming the swap counterparty is a corporation.

\vspace{6pt}
\noindent{\sffamily\bfseries\fontsize{8.6}{10}\selectfont\color{forest}Solution}\par
\[
(\$175\text{ million}+\$3.375\text{ million})\times0.08=\$14.27\text{ million}
\]
\end{orexambox}

% ---------------------------------------------------------------------
\lohead{60}{c}{Describe measures introduced in the 1995 and 1996 amendments,
including guidelines for netting of credit exposures and methods for calculating
market risk capital for assets in the trading book.}

\orsub{The 1995 amendment: netting}

By 1995, the 1987 crash and the popularity of VaR highlighted market risk gaps in
Basel I. Basel I discouraged hedging prior to 1995, so netting was allowed in 1995.

\begin{itemize}
\item \textbf{Netting:} allowed bilateral netting of credit exposures, reducing the
credit equivalent amount (CEA) needed for regulatory capital.
\item \textbf{Net Replacement Ratio (NRR):} introduced to measure the impact of
netting by comparing current exposure with and without netting. This ratio helped
banks reduce their risk-weighted assets (RWA) by showing a lower net exposure.
\end{itemize}

\begin{orkeybox}
\[
\mathrm{CEA}=\max\!\left(\sum_{i=1}^{N}V_i,\;0\right)
 +\sum_{j}\bigl[0.4\times D_j + 0.6\times D_j \times \mathrm{NRR}\bigr]
\]
\end{orkeybox}

\orsub{The 1996 amendment: market risk and trading activities}

The 1995 amendment didn't address market risk. The 1996 amendment required banks
to:

\begin{lettered}
\item Measure market risks from trading activities.
\item Hold capital to cover those risks.
\end{lettered}

\orsubb{Two methods to calculate market risk in the amendment}

\begin{enumerate}
\item[\textcolor{rust}{\sffamily\bfseries 1)}] \textbf{Standardized approach.}
Simpler method with pre-set capital charges for each trading book item, used by
less sophisticated banks. Ignores correlations between instruments.
\item[\textcolor{rust}{\sffamily\bfseries 2)}] \textbf{Internal model-based
approach.}
  \begin{lettered}
  \item Allows banks to use their own risk measurement models, typically
  calculating Value at Risk (VaR) over a 10-day period with a 99\% confidence
  level.
  \item Generally results in lower capital charges due to recognition of
  diversification benefits.
  \item Covers broad market risks including interest rates, currency rates, stock
  indices, and commodities.
  \end{lettered}
\end{enumerate}

\begin{orkeybox}
{\sffamily\bfseries\fontsize{8.6}{10}\selectfont\color{navy}Total capital until
1996}\ {\fontsize{8.8}{11}\selectfont\color{midgray}(only market and credit risk
capital was required)}
\[
\text{total capital}=0.08\times(\text{credit risk RWA}+\text{market risk RWA})
\]
where
\[
\text{market RWA}=12.5\times\bigl[\max(\mathrm{VaR}_{t-1},\ m\times\mathrm{VaR}_{avg})\bigr]
\]
\[
\text{credit RWA}=\textstyle\sum(\text{RWA on-balance sheet})+\sum(\text{RWA off-balance sheet})
\]
\end{orkeybox}

\begin{orexambox}{Example: market risk capital charge}
A bank calculates the previous day's market risk VaR as \$10 million. The average
VaR over the preceding 60 trading days is \$8 million. Assuming a multiplicative
factor of three, \textbf{calculate} the market risk capital charge.

\vspace{6pt}
\noindent{\sffamily\bfseries\fontsize{8.6}{10}\selectfont\color{forest}Answer}\par
\[
\text{market risk capital charge}=0.08\times\bigl[12.5\times(3\times\$8\text{ million})\bigr]=\$24\text{ million}
\]
\end{orexambox}

\begin{ornotebox}{Backtesting requirement}
Banks must backtest their one-day 99\% VaR models over the previous 250 trading
days to ensure accuracy, adjusting the multiplicative factor based on the number of
exceptions (losses exceeding VaR predictions).
\end{ornotebox}

% ---------------------------------------------------------------------
\lohead{60}{d}{Describe changes to the Basel regulations made as part of Basel II,
including the three pillars.}

\begin{ornotebox}{Basel I limitations}
Ignored borrower creditworthiness (all loans had 100\% risk weight), and neglected
diversification benefits (no consideration of default correlation).
\end{ornotebox}

\orsub{Basel II's improvements: the three pillars}

\begin{enumerate}
\item[\textcolor{rust}{\sffamily\bfseries 1)}] \textbf{Pillar 1: minimum capital
requirements.} Risk-weighted assets (RWA) based on borrower creditworthiness.
Capital charges for market risk remained the same as Basel I. Introduced capital
charges for operational risk. Maintained the 8\% total capital ratio requirement.
\item[\textcolor{rust}{\sffamily\bfseries 2)}] \textbf{Pillar 2: supervisory
review.} Encourages supervisors to evaluate bank risk management functions and
consider local conditions when implementing Basel II. Introduced the Internal
Capital Adequacy Assessment Process (ICAAP) for banks.
\item[\textcolor{rust}{\sffamily\bfseries 3)}] \textbf{Pillar 3: market
discipline.} Promotes transparency by requiring banks to disclose more risk
information. Disclosures include: (a) bank structure and capital levels, (b) risk
exposures and risk measures, and (c) capital allocation for different risks.
\end{enumerate}

% ---------------------------------------------------------------------
\lohead{60}{e}{Compare the standardized internal ratings-based (IRB) approach, the
foundation IRB approach, and the advanced IRB approach for the calculation of
credit risk capital under Basel II.}

\lohead{60}{f}{Calculate credit risk capital under Basel II utilizing the IRB
approach.}

\orsub{Credit risk capital requirements, Basel II}

\orsubb{1) The standardized approach}

\begin{lettered}
\item OECD status is no longer valid.
\item Risk weights assigned based on credit ratings.
\item Allows for lower risk weights for exposures with short maturities under
certain conditions.
\end{lettered}

\begin{orexambox}{Example: collateral adjustment, simple approach}
Blue Star Bank has a \$100 million exposure to Monarch, Inc. The exposure is
secured by \$80 million of collateral consisting of AAA-rated bonds. Monarch has a
credit rating of B. The collateral risk weight is 20\% and the counterparty risk
weight is 150\%. Using the simple approach, \textbf{calculate} the risk-weighted
assets.

\vspace{6pt}
\noindent{\sffamily\bfseries\fontsize{8.6}{10}\selectfont\color{forest}Solution}\par
\[
(0.2\times80)+(1.5\times20)=\$46\text{ million risk-weighted assets}
\]
\end{orexambox}

\begin{orexambox}{Example: comprehensive approach}
Blue Star Bank assumes an adjustment to the exposure in the previous example of
$+15\%$ to allow for possible increases in the exposures. The bank also allows for
a $-20\%$ change in the value of the collateral. \textbf{Calculate} the new
exposure using the comprehensive approach.

\vspace{6pt}
\noindent{\sffamily\bfseries\fontsize{8.6}{10}\selectfont\color{forest}Solution}\par
\[
(1.15\times100)-(0.8\times80)=\$51\text{ million exposure}
\qquad\Longrightarrow\qquad
1.5\times51=\$76.5\text{ million RWA}
\]
\end{orexambox}

\orsubb{2) The Internal Ratings-Based (IRB) approach}

\begin{lettered}
\item Requires stronger risk management capabilities.
\item Capital based on a Value at Risk (VaR) calculation over one year at a 99.9\%
confidence level.
\item Focuses on unexpected losses (UL), expecting banks to cover expected losses
through pricing.
\end{lettered}

\begin{orkeybox}
{\sffamily\bfseries\fontsize{8.6}{10}\selectfont\color{navy}Formula 1}
\[
\text{capital}=\sum_i\bigl[\mathrm{EAD}_i\times \mathrm{LGD}_i\times \mathrm{DR}99.9_i\bigr]
-\sum_i\bigl[\mathrm{EAD}_i\times \mathrm{LGD}_i\times \mathrm{PD}_i\bigr]
\]
\[
\text{capital}=\mathrm{VaR}_{99.9\%,\,1\text{-year}}-\mathrm{EL}
\]
\end{orkeybox}

\begin{ornotebox}{Remember}
If the correlation is 0 among borrowers, then $\mathrm{DR}99.9=\mathrm{PD}$ in the
formula above.
\end{ornotebox}

\begin{orkeybox}
{\sffamily\bfseries\fontsize{8.6}{10}\selectfont\color{navy}Formula 2}
\[
\text{capital}=\mathrm{EAD}\times \mathrm{LGD}\times(\mathrm{DR}99.9-\mathrm{PD})\times \mathrm{MA}
\]
where
\[
\mathrm{MA}=\text{maturity adjustment}=\frac{1+(M-2.5)\times b}{1-1.5\times b}
\]
\[
M=\text{maturity of the exposure}\qquad
b=\bigl[0.11852-0.05478\times\ln(\mathrm{PD})\bigr]^{2}
\]
\end{orkeybox}

\begin{orkeybox}
{\sffamily\bfseries\fontsize{8.6}{10}\selectfont\color{navy}Formula 3}
\[
\mathrm{RWA}=12.5\times\bigl[\mathrm{EAD}\times \mathrm{LGD}\times(\mathrm{DR}99.9-\mathrm{PD})\times \mathrm{MA}\bigr]
\]
\vspace{-4pt}
{\fontsize{9}{11}\selectfont\color{rust} and then multiply RWA by 8\% to get
capital.}
\end{orkeybox}

\orsubb{Comparing the two IRB approaches}

\noindent\begin{minipage}{\textwidth}
\renewcommand{\arraystretch}{1.32}
\begin{tabularx}{\textwidth}{@{}L{0.13\textwidth}YY@{}}
\rowcolor{navy} \thd{} & \thd{Foundation IRB} & \thd{Advanced IRB}\\
{\tabfont \tsub{PD}} & {\tabfont Bank-supplied estimates with a 0.03\% floor} & {\tabfont Bank-supplied, customisable, with a 0.03\% floor}\\
\rowcolor{paleblue}
{\tabfont \tsub{LGD}} & {\tabfont Set by regulators (45\% senior, 75\% subordinated); reduced by collateral} & {\tabfont Bank-supplied, based on collateral and seniority}\\
{\tabfont \tsub{EAD}} & {\tabfont Similar to Basel I's CEA with netting} & {\tabfont Bank uses own estimates with supervisory approval}\\
\rowcolor{paleblue}
{\tabfont \tsub{Maturity}} & {\tabfont Fixed at 2.5} & {\tabfont Bank-supplied, reflects actual maturity}\\
\end{tabularx}
\end{minipage}
\orfigcap{Foundation versus Advanced IRB, parameter by parameter.}

\orsubb{IRB approach for retail exposures (combination)}

\begin{lettered}
\item Merges Foundation and Advanced IRB approaches.
\item Bank estimates PD, EAD, and LGD.
\item No maturity adjustment.
\item Capital requirement based on EAD, LGD, and the difference between DR99.9 and
PD.
\item Lower correlations assumed for retail exposures.
\end{lettered}

\begin{orexambox}{Example: RWA under the IRB approach}
Assume Blue Star Bank has a \$150 million loan to an A-rated corporation. The PD is
0.1\% and the LGD is 50\%. The DR99.9 is 3.4\%. The average maturity of the loan is
2.5 years. \textbf{Calculate} the RWA using the IRB approach and compare it to the
RWA under Basel I.

\vspace{6pt}
\noindent{\sffamily\bfseries\fontsize{8.6}{10}\selectfont\color{forest}Solution}\par
\[
b=\bigl[0.11852-0.05478\times\ln(0.001)\bigr]^{2}=0.247
\]
\[
\mathrm{MA}=\frac{1}{1-(1.5\times0.247)}=1.59
\]
\[
\text{risk-weighted assets}=12.5\times150\times0.5\times(0.034-0.001)\times1.59
=\$49.19\text{ million}
\]
Under Basel I, the RWA for corporate loans was 100\%, or \$150 million in this
case. Thus, the IRB approach lowers the RWA for higher-rated corporate loans, in
this case from \$150 million to \$49.19 million.
\end{orexambox}

% ---------------------------------------------------------------------
\lohead{60}{g}{Compare the basic indicator approach, the standardized approach, and
the advanced measurement approach for the calculation of operational risk capital
under Basel II.}

\orsub{Operational risk capital requirements}

\noindent\begin{minipage}{\textwidth}
\renewcommand{\arraystretch}{1.32}
\begin{tabularx}{\textwidth}{@{}L{0.26\textwidth}Y@{}}
\rowcolor{navy} \thd{Approach} & \thd{How capital is set}\\
{\tabfont \textbf{1) Basic Indicator Approach (BIA)}} &
{\tabfont Simplest method, used by banks with less advanced risk management.
Capital is set at 15\% of the three-year average annual gross income, excluding
years with negative income.}\\
\rowcolor{paleblue}
{\tabfont \textbf{2) Standardized Approach}} &
{\tabfont Builds on the BIA by applying different multipliers to the gross income
from various business lines (e.g., retail banking might have a 12\% multiplier,
commercial banking 15\%).}\\
{\tabfont \textbf{3) Advanced Measurement Approach (AMA)}} &
{\tabfont Most sophisticated, similar to the IRB approach for credit risk. Uses the
bank's internal models to estimate operational risk as Value at Risk (VaR) over a
one-year period at a 99.9\% confidence level.}\\
\end{tabularx}
\end{minipage}
\orfigcap{Three routes to operational risk capital under Basel II.}

% ---------------------------------------------------------------------
\lohead{60}{h}{Summarize elements of the Solvency II capital framework for
insurance companies.}

\orsub{Capital requirements for insurance}

\begin{lettered}
\item \textbf{Minimum Capital Requirement (MCR):} base level of capital.
\item \textbf{Solvency Capital Requirement (SCR):} MCR plus capital buffers.
\end{lettered}

\orsubb{Calculating SCR: two approaches}

\begin{enumerate}
\item[\textcolor{rust}{\sffamily\bfseries 1)}] \textbf{Standardized approach.}
Simpler method for less sophisticated firms, similar to Basel II's standardized
approach.
\item[\textcolor{rust}{\sffamily\bfseries 2)}] \textbf{Internal models approach.}
More complex, firm-specific risk model similar to Basel II's IRB approach. Uses VaR
at a 99.5\% confidence level for one year.
\end{enumerate}
